\begin{textsample}{POS Dim 1 – 2013 – Score 21.00 – 2013/t001929.txt}  \label{ex:f1_pos_043}
Technology can change our understanding of nature.
Take for example \textbf{the} case of \textbf{lions}.
For centuries, it ' s been said that female \textbf{lions} do all of \textbf{the} hunting out in \textbf{the} open savanna, and male \textbf{lions} do nothing until it ' s time for dinner.
You ' ve heard this too, I can tell.
Well recently, I led an airborne mapping campaign in \textbf{the} Kruger National Park in South Africa.
Our colleagues put GPS tracking collars on male and female \textbf{lions}, and we mapped their hunting behavior from \textbf{the} air. \textbf{The} lower \textbf{left} shows a \textbf{lion} sizing up a herd of impala for a kill, and \textbf{the} right shows what I call \textbf{the} \textbf{lion} viewshed.
That ' s how far \textbf{the} \textbf{lion} can see in all directions until his or her view is obstructed by vegetation.
And what we found is that male \textbf{lions} are not \textbf{the} lazy hunters we thought them to be.
They just use a different strategy.
Whereas \textbf{the} female \textbf{lions} hunt out in \textbf{the} open savanna over long distances, usually during \textbf{the} day, male \textbf{lions} use an ambush strategy in dense vegetation, and often at night.
This video shows \textbf{the} actual hunting viewsheds of male \textbf{lions} on \textbf{the} \textbf{left} and females on \textbf{the} right.
Red and darker colors show more dense vegetation, and \textbf{the} white are wide open spaces.
And this is \textbf{the} viewshed right literally at \textbf{the} eye level of hunting male and female \textbf{lions}.
All of a sudden, you get a very clear understanding of \textbf{the} very spooky conditions under which male \textbf{lions} do their hunting.
I bring up this example to begin, because it emphasizes how little we know about nature.
There ' s been a huge amount of work done so far to try to slow down our losses of tropical forests, and we are losing our forests at a rapid rate, as shown in red on \textbf{the} \textbf{slide}.
I find it ironic that we ' re doing so much, yet these areas are fairly unknown to science.
So how can we save what we don ' t understand?
Now I ' m a global ecologist and an Earth explorer with a background in physics and \textbf{chemistry} and \textbf{biology} and a lot of other boring subjects, but above all, I ' m obsessed with what we don ' t know about our \textbf{planet}.
So I created this, \textbf{the} Carnegie Airborne Observatory, or CAO.
It may look like a plane with a fancy paint job, but I packed it with over 1, 000 kilos of high-tech sensors, \textbf{computers}, and a very motivated staff of Earth scientists and pilots.
Two of our instruments are very unique : one is called an imaging spectrometer that can actually measure \textbf{the} chemical composition of plants as we fly over them.
Another one is a set of lasers, very high-powered lasers, that fire out of \textbf{the} bottom of \textbf{the} plane, sweeping across \textbf{the} ecosystem and measuring it at nearly 500, 000 times per second in high-resolution 3D.
Here ' s an image of \textbf{the} Golden Gate Bridge in San Francisco, not far from where I live.
Although we flew straight over this bridge, we imaged it in 3D, captured its color in just a few seconds.
But \textbf{the} real power of \textbf{the} CAO is its ability to capture \textbf{the} actual building blocks of ecosystems.
This is a small town in \textbf{the} Amazon, imaged with \textbf{the} CAO.
We can slice through our data and see, for example, \textbf{the} 3D structure of \textbf{the} vegetation and \textbf{the} buildings, or we can use \textbf{the} chemical information to actually figure out how fast \textbf{the} plants are growing as we fly over them. \textbf{The} hottest pinks are \textbf{the} fastest-growing plants.
And we can see biodiversity in ways that you never could have imagined.
This is what a rainforest might look like as you fly over it in a hot air balloon.
This is how we see a rainforest, in kaleidoscopic color that tells us that there are many species living with one another.
But you have to remember that these trees are literally bigger than \textbf{whales}, and what that means is that they ' re impossible to understand just by walking on \textbf{the} ground below them.
So our imagery is 3D, it ' s chemical, it ' s biological, and this tells us not only \textbf{the} species that are living in \textbf{the} canopy, but it tells us a lot of information about \textbf{the} rest of \textbf{the} species that occupy \textbf{the} rainforest.
Now I created \textbf{the} CAO in order to answer questions that have proven extremely challenging to answer from any other vantage point, such as from \textbf{the} ground, or from satellite sensors.
I want to share three of those questions with you today. \textbf{The} first questions is, how do we manage our \textbf{carbon} reserves in tropical forests?
Tropical forests contain a huge amount of \textbf{carbon} in \textbf{the} trees, and we need to keep that \textbf{carbon} in those forests if we ' re going to avoid any further global warming.
Unfortunately, global \textbf{carbon} \textbf{emissions} from deforestation now equals \textbf{the} global \textbf{transportation} sector.
That ' s all ships, \textbf{airplanes}, trains and automobiles combined.
So it ' s understandable that policy negotiators have been working hard to reduce deforestation, but they ' re doing it on landscapes that are hardly known to science.
If you don ' t know where \textbf{the} \textbf{carbon} is exactly, in detail, how can you know what you ' re losing?
Basically, we need a high-tech accounting system.
With our system, we ' re able to see \textbf{the} \textbf{carbon} stocks of tropical forests in utter detail. \textbf{The} red shows, obviously, closed- canopy tropical forest, and then you see \textbf{the} cookie cutting, or \textbf{the} cutting of \textbf{the} forest in \textbf{yellows} and greens.
It ' s like cutting a cake except this cake is about \textbf{whale} deep.
And yet, we can zoom in and see \textbf{the} forest and \textbf{the} trees at \textbf{the} same time.
And what ' s amazing is, even though we flew very high above this forest, later on in analysis, we can go in and actually experience \textbf{the} treetrops, leaf by leaf, branch by branch, just as \textbf{the} other species that live in this forest experience it along with \textbf{the} trees themselves.
We ' ve been using \textbf{the} technology to explore and to actually put out \textbf{the} first \textbf{carbon} geographies in high resolution in faraway places like \textbf{the} Amazon Basin and not-so-faraway places like \textbf{the} United States and Central America.
What I ' m going to do is I ' m going to take you on a high-resolution, first-time tour of \textbf{the} \textbf{carbon} landscapes of Peru and then Panama. \textbf{The} colors are going to be going from red to \textbf{blue}.
Red is extremely high \textbf{carbon} stocks, your largest cathedral forests you can imagine, and \textbf{blue} are very low \textbf{carbon} stocks.
And let me tell you, Peru alone is an amazing place, totally unknown in terms of its \textbf{carbon} geography until today.
We can fly to this area in northern Peru and see super high \textbf{carbon} stocks in red, and \textbf{the} Amazon River and floodplain cutting right through it.
We can go to an area of utter devastation caused by deforestation in \textbf{blue}, and \textbf{the} virus of deforestation spreading out in orange.
We can also fly to \textbf{the} southern Andes to see \textbf{the} tree line and see exactly how \textbf{the} \textbf{carbon} geography ends as we go up into \textbf{the} mountain system.
And we can go to \textbf{the} biggest swamp in \textbf{the} western Amazon.
It ' s a watery dreamworld akin to Jim Cameron ' s"Avatar."We can go to one of \textbf{the} smallest tropical countries, Panama, and see also a huge range of \textbf{carbon} variation, from high in red to low in \textbf{blue}.
Unfortunately, most of \textbf{the} \textbf{carbon} is lost in \textbf{the} lowlands, but what you see that ' s left, in terms of high \textbf{carbon} stocks in greens and reds, is \textbf{the} stuff that ' s up in \textbf{the} mountains.
One interesting exception to this is right in \textbf{the} middle of your screen.
You ' re seeing \textbf{the} buffer zone around \textbf{the} Panama Canal.
That ' s in \textbf{the} reds and \textbf{yellows}. \textbf{The} canal authorities are using force to protect their watershed and global commerce.
This kind of \textbf{carbon} mapping has transformed conservation and resource policy development.
It ' s really advancing our ability to save forests and to curb climate change.
My second question : How do we prepare for climate change in a place like \textbf{the} Amazon rainforest?
Let me tell you, I spend a lot of time in these places, and we ' re seeing \textbf{the} climate changing already.
Temperatures are increasing, and what ' s really happening is we ' re getting a lot of \textbf{droughts}, recurring \textbf{droughts}. \textbf{The} 2010 mega-drought is shown here with red showing an area about \textbf{the} size of Western Europe. \textbf{The} Amazon was so dry in 2010 that even \textbf{the} main stem of \textbf{the} Amazon river itself dried up partially, as you see in \textbf{the} photo in \textbf{the} lower portion of \textbf{the} \textbf{slide}.
What we found is that in very remote areas, these \textbf{droughts} are having a big negative impact on tropical forests.
For example, these are all of \textbf{the} dead trees in red that suffered mortality following \textbf{the} 2010 \textbf{drought}.
This area happens to be on \textbf{the} border of Peru and Brazil, totally unexplored, almost totally unknown scientifically.
So what we think, as Earth scientists, is species are going to have to migrate with climate change from \textbf{the} east in Brazil all \textbf{the} way west into \textbf{the} Andes and up into \textbf{the} mountains in order to minimize their exposure to climate change.
One of \textbf{the} problems with this is that humans are taking apart \textbf{the} western Amazon as we speak.
Look at this 100-square-kilometer gash in \textbf{the} forest created by gold miners.
You see \textbf{the} forest in green in 3D, and you see \textbf{the} effects of gold mining down below \textbf{the} soil \textbf{surface}.
Species have nowhere to migrate in a system like this, obviously.
If you haven ' t been to \textbf{the} Amazon, you should go.
It ' s an amazing experience every time, no matter where you go.
You ' re going to probably see it this way, on a river.
But what happens is a lot of times \textbf{the} rivers hide what ' s really going on back in \textbf{the} forest itself.
We flew over this same river, imaged \textbf{the} system in 3D. \textbf{The} forest is on \textbf{the} \textbf{left}.
And then we can digitally remove \textbf{the} forest and see what ' s going on below \textbf{the} canopy.
And in this case, we found gold mining activity, all of it illegal, set back away from \textbf{the} river ' s edge, as you ' ll see in those strange pockmarks coming up on your screen on \textbf{the} right.
Don ' t worry, we ' re working with \textbf{the} authorities to deal with this and many, many other problems in \textbf{the} region.
So in order to put together a conservation plan for these unique, important corridors like \textbf{the} western Amazon and \textbf{the} Andes Amazon corridor, we have to start making geographically explicit plans now.
How do we do that if we don ' t know \textbf{the} geography of biodiversity in \textbf{the} region, if it ' s so unknown to science?
So what we ' ve been doing is using \textbf{the} laser-guided spectroscopy from \textbf{the} CAO to map for \textbf{the} first time \textbf{the} biodiversity of \textbf{the} Amazon rainforest.
Here you see actual data showing different species in different colors.
Reds are one type of species, \textbf{blues} are another, and greens are yet another.
And when we take this together and scale up to \textbf{the} regional level, we get a completely new geography of biodiversity unknown prior to this work.
This tells us where \textbf{the} big biodiversity changes \textbf{occur} from \textbf{habitat} to \textbf{habitat}, and that ' s really important because it tells us a lot about where species may migrate to and migrate from as \textbf{the} climate shifts.
And this is \textbf{the} pivotal information that ' s needed by decision makers to develop protected areas in \textbf{the} context of their regional development plans.
And third and final question is, how do we manage biodiversity on a \textbf{planet} of protected ecosystems? \textbf{The} example I started out with about \textbf{lions} hunting, that was a study we did behind \textbf{the} fence line of a protected area in South Africa.
And \textbf{the} truth is, much of Africa ' s nature is going to persist into \textbf{the} future in protected areas like I show in \textbf{blue} on \textbf{the} screen.
This puts incredible pressure and responsibility on park management.
They need to do and make decisions that will benefit all of \textbf{the} species that they ' re protecting.
Some of their decisions have really big impacts.
For example, how much and where to use fire as a management tool?
Or, how to deal with a large species like \textbf{elephants}, which may, if their populations get too large, have a negative impact on \textbf{the} ecosystem and on other species.
And let me tell you, these types of dynamics really play out on \textbf{the} landscape.
In \textbf{the} foreground is an area with lots of fire and lots of \textbf{elephants} : wide open savanna in \textbf{blue}, and just a few trees.
As we cross this fence line, now we ' re getting into an area that has had protection from fire and zero \textbf{elephants} : dense vegetation, a radically different ecosystem.
And in a place like Kruger, \textbf{the} soaring \textbf{elephant} densities are a real problem.
I know it ' s a sensitive issue for many of you, and there are no easy answers with this.
But what ' s good is that \textbf{the} technology we ' ve developed and we ' re working with in South Africa, for example, is allowing us to map every single tree in \textbf{the} savanna, and then through repeat flights we ' re able to see which trees are being pushed over by \textbf{elephants}, in \textbf{the} red as you see on \textbf{the} screen, and how much that ' s happening in different types of landscapes in \textbf{the} savanna.
That ' s giving park managers a very first opportunity to use tactical management strategies that are more nuanced and don ' t lead to those extremes that I just showed you.
So really, \textbf{the} way we ' re looking at protected areas nowadays is to think of it as tending to a circle of life, where we have fire management, \textbf{elephant} management, those impacts on \textbf{the} structure of \textbf{the} ecosystem, and then those impacts affecting everything from insects up to apex predators like \textbf{lions}.
Going forward, I plan to greatly expand \textbf{the} airborne observatory.
I ' m hoping to actually put \textbf{the} technology into orbit so we can manage \textbf{the} entire \textbf{planet} with technologies like this.
Until then, you ' re going to find me flying in some remote place that you ' ve never heard of.
I just want to end by saying that technology is absolutely critical to managing our \textbf{planet}, but even more important is \textbf{the} understanding and wisdom to apply it.
Thank you.

% matched lemmas: airplane, biology, blue, carbon, chemistry, computer, drought, elephant, emission, habitat, left, lion, occur, planet, slide, surface, the, transportation, whale, yellow
\end{textsample}
